% !TeX spellcheck = en_US

\section{\textsf{S-CORE}}
\label{section:S-CORE}
We introduce \textsf{S-CORE} as an extension of \textsf{SRL} \cite{Matos:TCS03}, a reversible computational model that exclusively represents total injective functions. 
Two key features distinguish \SCORE from \SRL:
\begin{enumerate}
    \item \SCORE includes the primitives \lstiC{PUSH} and \lstiC{POP} for stack management;
    \item Each variable \lstiC{x} in \SCORE is represented as a triple consisting of:
    \begin{itemize}
        \item The first component, storing the current value of \lstiC{x};
        \item The second component, maintaining the history of values that \lstiC{x} has assumed;
        \item The third component, enabling \lstiC{PUSH} to function as a successor operation and \lstiC{POP} as a predecessor operation, ensuring their mutual reversibility.
    \end{itemize}
\end{enumerate}
\noindent
Before proceeding, let us recall the main features of \SRL.

%%-----------------
\subsection{\SRL in a nutshell.}
\label{subsection:SRL in a nutshell}
\textsf{SRL} is a reversible imperative computational model 
in which every term \lstiC{P} syntactically determines its inverse \lstiC{-P}. We will illustrate this aspect technically when introducing \SCORE, the subject of this work.
The variables of \SRL range over $\ZZ$. The instruction ``\lstiC{INC x}'' increments the value of \lstiC{x} by one, whereas ``\lstiC{DEC x}'' decrements it by one. Moreover,
given that $v$ represents the value of \lstiC{x}, the loop construct ``\lstiC{FOR x P}'' iterates the body \lstiC{P} exactly $|v|$ times, following these rules:
\begin{itemize}
    \item If $v > 0$, then \lstiC{P} is \emph{iterated forward}, meaning each iteration executes \lstiC{P}.
    \item If $v < 0$, then the body is \emph{iterated backward}, meaning each iteration executes the \emph{inverse} \lstiC{-P} of \lstiC{P}.
    \item If $v = 0$, then \lstiC{P} is skipped entirely.
\end{itemize}



%%-----------------
\subsection{\SCORE syntax}
\label{subsection:SCORE: syntax and first notions}
\begin{definition}[\textsf{S-CORE} syntax]
\label{definition:S-CORE syntax}
Let $V$ be a set of variable names \lstiC{x}, \lstiC{y}, \lstiC{z}, \ldots \enspace. Let $\lstiC{x} \in V$. \SCORE terms belong to the language generated by the grammar:
\begin{align*}
	\lstiC{P} & ::= \lstiC{SKIP} \mid \lstiC{INC x} \mid \lstiC{DEC x}
	\mid \lstiC{PUSH x} \mid \lstiC{POP x} \mid \lstiC{P;P}  \mid
    \lstiC{FOR x \{P\}}
    \enspace ,
\end{align*}
under the proviso that the \emph{leading variable} \lstiC{x} in ``\lstiC{FOR x \{P\}}'' \emph{cannot} occur in \lstiC{P}, that is \lstiC{P} never contains neither ``\lstiC{INC x}'' nor ``\lstiC{DEC x}''. Writing $\lstiC{P} \in \SCORE$ means that \lstiC{P} is an \SCORE term.
\end{definition}

\begin{remark}
\SRL contains all and only the \SCORE terms which are free of ``\lstiC{PUSH x}'' and ``\lstiC{POP x}'', for every variable \lstiC{x}.
\end{remark}

\begin{definition}[Inverse of a term in \SCORE]
\label{definition:Inverse -P of any SCORE term P}
The function \invSC{(.)}$:\SCORE$ $\to \SCORE$ inverts every $\lstiC{P}\in\SCORE$ by structural induction on \lstiC{P}:
\begin{alignat*}{3}
\invSC{(INC x)} &= \lstiC{DEC x}
&\quad
\invSC{SKIP}    &= \lstiC{SKIP}
&\quad
\invSC{(DEC x)} &= \lstiC{INC x}
\\
\invSC{(PUSH x)} & = \lstiC{POP x}
&& &
\invSC{(POP x)}  & = \lstiC{PUSH x}
\\
\invSC{(P;Q)} & = \invSC{Q}\lstiC{;}\invSC{P}
&& &
\invSC{(FOR x \{P\})} & = \lstiC{FOR x \{}\invSC{P}\lstiC{\}}
\enspace .
\end{alignat*}
\end{definition}

\begin{remark}
It should be obvious to see that the function \invSC{(.)} is \emph{self-dual}, i.e. $\lstiC{-(-(P))} = \lstiC{P}$, for every \SCORE program \lstiC{P}. 
\end{remark}

%%%%%%%%%%%%%%%%%%%%
\subsection{Our plan towards the final operational semantics \rsemantics of \SCORE}

Clearly, in a classical computational model, \lstiC{PUSH} and \lstiC{POP} are not a pair of bijections because it is meaningless to \lstiC{POP} elements from an empty stack. We need to extend the structure of the domain in which we interpret \lstiC{PUSH} and \lstiC{POP} to eventually obtain that both ``\lstiC{-(PUSH)} = \lstiC{POP}'' and ``\lstiC{-(POP)} = \lstiC{PUSH}'' also hold at the operational semantics level.

We think it is useful to follow a step-wise extension of the domain to operationally interpret \lstiC{PUSH} and \lstiC{POP}, to show a concrete example of how to inject non reversible functions into reversible ones.

Our step-wise approach is realized by defining three big-step operational semantics -- \nsemantics, \asemantics, and \rsemantics\xspace -- and the corresponding domain of interpretation they rely on, following two general guidelines:
\begin{enumerate}
    \item Each semantics is based on its own notion of \emph{state} $\sigma$, mapping every variable to an appropriate current value.
    \item Each semantics is formalized as a relation with judgments of the form
    $\langle \lstiC{P}, \sigma \rangle \Downarrow \tau$
    which we interpret as 
        ``For every $\lstiC{P} \in \SCORE$, and for every state $\sigma$ and $\tau$, executing \lstiC{P} from $\sigma$ results in $\tau$.'' 
\end{enumerate}\noindent


